% !TeX program = xelatex
% ============================================================================
%  第 2 课　从余数到模运算：单位元与逆元
% ----------------------------------------------------------------------------
%  编译        XeLaTeX，编译两遍
%  样式        ../../00-共享/preamble.tex
%  说明        本节配发一页单页（见 04-单页/）
%  图片        圆盘类插图一律 \imgph 占位，不要在这里手画 tikz
% ============================================================================

\documentclass[10pt]{beamer}
% ============================================================================
%  讲义样式文件 —— 字体 / 配色 / 校徽栏 / 顶部导航栏 / 页脚 / 常用方框
%  所有个性化设置都集中在这里，改一次即可用于所有讲义
%
%  用法：主文件里写
%      \documentclass[10pt]{beamer}
%      \input{preamble.tex}
%  依赖：同目录下的 logo_white.png（顶部校徽）
%  编译：必须用 XeLaTeX，而且必须编译两遍
%        （顶部导航栏的节名与小圆点记录在 .nav 文件里，只编一遍不会更新）
% ============================================================================
\usepackage{ctex}  % 支持中文
\usepackage{amsmath, amssymb}  % 数学公式支持
\usepackage{fontspec}  % 字体设置
\usepackage{unicode-math}  % 支持Unicode数学字体
\usepackage{xcolor}  % 颜色
\usepackage{tikz}  % 用于绘制页脚色条
\usepackage{booktabs}  % 表格

% 图片搜索路径 ------------------------------------------------------------------
% 校本课的资料分目录存放，本文件可能被 02-课件/第XX课/ 、03-讲义/ 下的主文件调用，
% 这里把几种可能的相对位置都列上，logo_white.png 放在 00-共享/ 即可被找到。
% 最后一项 00-共享/图片/ 是给 \imgfig 用的插图目录（图与课件不在同一层，靠它找到）。
\graphicspath{{./}{./图片/}{../../00-共享/}{../00-共享/}{00-共享/}{../../00-共享/图片/}{../00-共享/图片/}{00-共享/图片/}}

% 设置字体
\setCJKmainfont{SimSun}  % 设置中文字体为宋体
\setmainfont{Palatino Linotype}  % 设置英文字体为Palatino Linotype
\setmathfont{Latin Modern Math}  % 设置数学字体为Latin Modern Math
% 注意：beamer 默认把正文字体族设成无衬线（sans），Windows 下对应的中文字体是
% 微软雅黑，所以上面这行 SimSun 目前对正文不起作用。想让正文变成宋体，取消
% 下面这行的注释即可（标题栏、页脚、方框标题不受影响）。
% \renewcommand{\familydefault}{\rmdefault}

% 定义颜色 - 均为绿色系
\definecolor{maingreen}{RGB}{34,139,34}  % 森林绿（导航栏底色）
\definecolor{lightgreen}{RGB}{240,255,240}  % 蜜瓜白（帧标题栏底色）
\definecolor{accentgreen}{RGB}{46,139,87}  % 海绿（强调色）
\definecolor{textgreen}{RGB}{0,128,0}  % 纯绿（标题文字）
\definecolor{darkgreen}{RGB}{20,90,20}  % 深绿（顶部校徽栏底色）

% 设置Beamer主题
\usetheme{default}
\usecolortheme{default}

% 自定义主题
\setbeamercolor{background canvas}{bg=white}
\setbeamercolor{frametitle}{fg=textgreen, bg=lightgreen}
\setbeamercolor{title}{fg=textgreen}
\setbeamercolor{subtitle}{fg=textgreen}
\setbeamercolor{section in toc}{fg=accentgreen}
\setbeamercolor{normal text}{fg=black}
\setbeamercolor{structure}{fg=accentgreen}
\setbeamercolor{itemize item}{fg=accentgreen}
\setbeamercolor{itemize subitem}{fg=accentgreen}
\setbeamercolor{itemize subsubitem}{fg=accentgreen}
\setbeamercolor{block title}{fg=white, bg=accentgreen}
\setbeamercolor{block body}{bg=lightgreen}
% example 环境（底层是 exampleblock）用的是独立的一组颜色名，default 配色主题下
% 它们的底色是空的，所以只设 block 的话“例”会变成一行没有框的标题，这里接过来
\setbeamercolor{block title example}{parent=block title}
\setbeamercolor{block body example}{parent=block body}

% ==================== 常用定理类方框（五类配色）====================
% 五类各自的标题条底色，正文底色由同色自动调浅（!8!white）
\definecolor{boxThm}{RGB}{20,90,20}    % 一、定理类：深绿
\definecolor{boxDef}{RGB}{46,139,87}   % 二、定义类：海绿
\definecolor{boxExm}{RGB}{34,139,34}   % 三、例题类：森林绿
\definecolor{boxPrf}{RGB}{85,107,47}   % 四、过程类：橄榄绿
\definecolor{boxNote}{RGB}{0,121,107}  % 五、注记类：青绿

% 方框本体：先套本类配色，再调用 beamer 的 block（改动只在环境内部生效）
\newcommand{\gbox}[2]{%
  \setbeamercolor{block title}{fg=white,bg=#1}%
  \setbeamercolor{block body}{bg=#1!8!white}%
  \begin{block}{#2}}
\newcommand{\gboxend}{\end{block}}
% 标题后的可选补充名：\begin{theorem}[勾股定理] → 定理（勾股定理）
\newcommand{\gsub}[1]{\ifx\relax#1\relax\else（#1）\fi}


% 一、定理类（定理/推论/性质/命题/引理）
\renewenvironment{theorem}[1][]{\gbox{boxThm}{定理\gsub{#1}}}{\gboxend}
\renewenvironment{corollary}[1][]{\gbox{boxThm}{推论\gsub{#1}}}{\gboxend}
\newenvironment{character}[1][]{\gbox{boxThm}{性质\gsub{#1}}}{\gboxend}
\newenvironment{proposition}[1][]{\gbox{boxThm}{命题\gsub{#1}}}{\gboxend}
\renewenvironment{lemma}[1][]{\gbox{boxThm}{引理\gsub{#1}}}{\gboxend}

% 二、定义类（定义/公理）
\renewenvironment{definition}[1][]{\gbox{boxDef}{定义\gsub{#1}}}{\gboxend}
\newenvironment{axiom}[1][]{\gbox{boxDef}{公理\gsub{#1}}}{\gboxend}

% 三、例题类（例/练习/习题）
\renewenvironment{example}[1][]{\gbox{boxExm}{例\gsub{#1}}}{\gboxend}
\newenvironment{exercise}[1][]{\gbox{boxExm}{练习\gsub{#1}}}{\gboxend}
\newenvironment{homework}[1][]{\gbox{boxExm}{习题\gsub{#1}}}{\gboxend}

% 四、过程类（证明/解/法1~法4）
\renewenvironment{proof}[1][]{\gbox{boxPrf}{证明\gsub{#1}}}{\qedsymbol\gboxend}
\newenvironment{jie}[1][]{\gbox{boxPrf}{解\gsub{#1}}}{\gboxend}
\newenvironment{way1}[1][]{\gbox{boxPrf}{法1\gsub{#1}}}{\gboxend}
\newenvironment{way2}[1][]{\gbox{boxPrf}{法2\gsub{#1}}}{\gboxend}
\newenvironment{way3}[1][]{\gbox{boxPrf}{法3\gsub{#1}}}{\gboxend}
\newenvironment{way4}[1][]{\gbox{boxPrf}{法4\gsub{#1}}}{\gboxend}

% 五、注记类（经验总结/注/值得注意的是/分析/提示）
\newenvironment{experience}[1][]{\gbox{boxNote}{经验总结\gsub{#1}}}{\gboxend}
\newenvironment{zhu}[1][]{\gbox{boxNote}{注\gsub{#1}}}{\gboxend}
\newenvironment{att}[1][]{\gbox{boxNote}{值得注意的是\gsub{#1}}}{\gboxend}
\newenvironment{analyse}[1][]{\gbox{boxNote}{分析\gsub{#1}}}{\gboxend}
\newenvironment{hint}[1][]{\gbox{boxNote}{提示\gsub{#1}}}{\gboxend}

% 顶部导航栏（miniframes 主题：显示各节名称 + 帧圆点）
% 节名与圆点记录在 .nav 文件中，改动章节后需要编译两遍才会同步
\useoutertheme[subsection=false]{miniframes}
\setbeamercolor{section in head/foot}{fg=yellow, bg=maingreen}
\setbeamercolor{mini frame}{fg=yellow, bg=maingreen}
\setbeamercolor{mini frame in current subsection}{fg=yellow, bg=accentgreen}

% 顶部校徽栏：位于导航栏上方，宽度与帧标题栏一致（整页宽）
% 校训用的楷体：ctex 的 \kaishu 在本机指向未安装的 Adobe 楷体，这里直接指定系统楷体
\setCJKfamilyfont{kaiti}{KaiTi}
\setbeamercolor{logo bar}{fg=white, bg=darkgreen}
\addtobeamertemplate{headline}{%
  \begin{beamercolorbox}[wd=\paperwidth,ht=0.84cm,dp=0.20cm,leftskip=0.25cm,rightskip=0.25cm]{logo bar}%
    % -0.035cm：色条内容默认坐在基线上，会比栏的中线略高，下移这一点使其上下居中
    \raisebox{-0.035cm}{\includegraphics[height=0.72cm]{logo_white.png}}%
    % 右侧校训：字体为上面定义的楷体，也可改为 \insertsectionhead（当前节名）、\inserttitle（章节名）等
    % 0.375cm = 校徽中线（0.36cm）加汉字基线偏移，并同样下移 0.035cm 与校徽对齐
    \hfill\raisebox{\dimexpr 0.375cm - 0.5\height - 0.5\depth\relax}%
      {\CJKfamily{kaiti}\fontsize{12pt}{24pt}\selectfont 自强不息，止于至善}%
  \end{beamercolorbox}%
}{}

% % 设置边框 - 使用TikZ绘制圆角
% \setbeamertemplate{frametitle}{
%   \vspace*{0.2cm}
%   \begin{tikzpicture}
%     % 背景矩形宽度限制为标题内容宽度
%     \fill[lightgreen, rounded corners=5pt] (0,0) rectangle (\textwidth+0.6cm, 0.8cm);
%     \node[anchor=west, text width=\textwidth, xshift=0.3cm, yshift=0.4cm] {\usebeamerfont{frametitle}\insertframetitle};
%   \end{tikzpicture}
% }

% 设置项目符号
\setbeamertemplate{itemize item}{\small\raise0.5pt\hbox{\textcolor{accentgreen}{$\bullet$}}}
\setbeamertemplate{itemize subitem}{\small\raise0.3pt\hbox{\textcolor{accentgreen}{$\circ$}}}

% 设置页脚 - 使用TikZ绘制一条整页宽的色条，左侧页码、右侧姓名
\setbeamertemplate{footline}{
  \begin{tikzpicture}
    \fill[maingreen] (-1,0) rectangle (\paperwidth, 0.4cm);
    \node[yellow,anchor=east, xshift=0cm, yshift=0.2cm] {\insertframenumber/\inserttotalframenumber};
    % 原点在纸张左边缘右侧 1cm 处（色条从 x=-1cm 画起），故 \paperwidth-2cm
    % 对应纸张右边缘内缩 1cm 的位置，与左侧页码左右对称
    \node[yellow,anchor=east, xshift=\paperwidth-2cm, yshift=0.2cm] {\insertauthor};
  \end{tikzpicture}
}

% 关闭 beamer 默认的导航符号（页码与姓名已在页脚中自行绘制）
\setbeamertemplate{navigation symbols}{}

% 设置标题页面
\setbeamertemplate{title page}{
  \vspace{0.5cm}
  \begin{center}
    \begin{minipage}[c][0.68\paperheight][c]{\textwidth}
      {\usebeamerfont{title}\usebeamercolor[fg]{title}\large\inserttitle\par}
      \vspace{1.5cm}
      {\usebeamerfont{subtitle}\usebeamercolor[fg]{subtitle}\Large\insertsubtitle\par}
      \vspace{1.5cm}
      {\usebeamerfont{author}\usebeamercolor[fg]{author}\insertauthor\par}
      \vspace{0.5cm}
      {\usebeamerfont{institute}\usebeamercolor[fg]{institute}\insertinstitute\par}
      \vspace{0.5cm}
      {\usebeamerfont{date}\usebeamercolor[fg]{date}\insertdate\par}
    \end{minipage}
  \end{center}
}

% 设置块环境
\setbeamertemplate{blocks}[rounded][shadow=false]

% 设置目录
\setbeamertemplate{section in toc}{\inserttocsectionnumber.~\inserttocsection}


% ----------------------------------------------------------------------------
%  图片占位符（临时用；有图以后把那行换成 \includegraphics 即可）
%      \imgph{说明文字}          默认高 3.6cm
%      \imgph[6cm]{说明文字}     指定高度
%  会画一个绿色虚线框，框里写说明，方便对着 PDF 找位置。
\newcommand{\imgph}[2][3.6cm]{%
  \begin{center}%
  \begin{tikzpicture}%
    \node[draw=accentgreen, densely dashed, line width=1pt, rounded corners=4pt,
          fill=lightgreen, inner sep=8pt, align=center, text width=0.72\textwidth,
          minimum height=#1, font=\small]
      {\textcolor{accentgreen}{\bfseries 【图片占位】}\\[4pt] #2};%
  \end{tikzpicture}%
  \end{center}%
}

% ----------------------------------------------------------------------------
%  插图（图已就位后用这个；图放在 00-共享\图片\ 下，由 \graphicspath 找到）
%      \imgfig{文件名}                 默认宽 0.72\textwidth
%      \imgfig[0.55\textwidth]{文件名}  指定宽度
%      \imgfig[8cm]{文件名}            也可以直接给尺寸
%  图源与编译方法见 00-共享\图片\README.md。每张图是一页 PDF，只含图、无页边。
\newcommand{\imgfig}[2][0.72\textwidth]{%
  \begin{center}%
    \includegraphics[width=#1]{#2}%
  \end{center}%
}

% ----------------------------------------------------------------------------
%  可用方框环境速查（五类配色，均不编号，均支持 \begin{xxx}[补充名]）
%    定理类（深绿）  theorem 定理 / corollary 推论 / character 性质 /
%                    proposition 命题 / lemma 引理
%    定义类（海绿）  definition 定义 / axiom 公理
%    例题类（森林绿）example 例 / exercise 练习 / homework 习题
%    过程类（橄榄绿）proof 证明 / jie 解 / way1 法1 / way2 法2 /
%                    way3 法3 / way4 法4
%    注记类（青绿）  experience 经验总结 / zhu 注 / att 值得注意的是 /
%                    analyse 分析 / hint 提示
%  另有 beamer 自带的普通方框：\begin{block}{标题} ... \end{block}
% ----------------------------------------------------------------------------


% ==================== 每份课件只改这一段 ====================
\title{从数数开始的抽象代数生活}
\subtitle{第 2 课　从余数到模运算：单位元与逆元}
\author{姜博瀚}
\institute{厦门大学附属实验中学}
\date{\today}
% ============================================================

\begin{document}

\begin{frame}
  \titlepage
\end{frame}

% ============================================================================
\section{取余数}

\begin{frame}{今天是星期三，100 天后是星期几？}
  \pause
  \begin{block}{先做一次除法}
    \[
      100 = 7 \times 14 + 2
    \]
  \end{block}
  \pause
  7 天一循环。14 个整星期可以整块扔掉，剩下 2 天。\pause 所以是\textbf{星期五}。
\end{frame}

\begin{frame}{带余除法，你们早就会了}
  \begin{block}{被除数 = 除数 × 商 + 余数}
    \[
      \underbrace{100}_{\text{被除数}} = \underbrace{7}_{\text{除数}} \times \underbrace{14}_{\text{商}} + \underbrace{2}_{\text{余数}}
    \]
  \end{block}
  \pause
  以前我们关心\textbf{商}。\pause 从今天起，我们只关心\textbf{余数}。
\end{frame}

% ============================================================================
\section{模记号}

\begin{frame}{给「余数」一个记号：mod}
  \[
    100 \equiv 2 \pmod 7
  \]
  读作「100 与 2 对模 7 同余」。
  \pause
  \begin{block}{它的意思就是一句话}
    \(100\) 除以 \(7\) 的余数是 \(2\)。
  \end{block}
\end{frame}

\begin{frame}{马上用一下}
  \begin{columns}[T]
    \begin{column}{0.48\textwidth}
      \begin{itemize}
        \item \(365 \equiv \text{?} \pmod 7\)
        \item \(1000 \equiv \text{?} \pmod 7\)
      \end{itemize}
    \end{column}
    \begin{column}{0.48\textwidth}
      \pause
      \begin{itemize}
        \item \(365 = 7\times 52 + 1 \Rightarrow 1\)
        \item \(1000 = 7\times 142 + 6 \Rightarrow 6\)
      \end{itemize}
    \end{column}
  \end{columns}
  \pause
  \vspace{0.6em}
  这两个余数，就是「过 365 天是星期几」和「过 1000 天是星期几」。
\end{frame}

% ============================================================================
\section{模 12 的加法}

\begin{frame}{把 12 个数摆成一个圈}
  \imgph[4.2cm]{12 格圆盘：\(0,1,2,\dots,11\) 顺时针排一圈\\（这张图全课反复用，建议做成一页单页）}
  \begin{center}
    模 \(12\) 的世界：\(0,1,2,\dots,11\)，走满 12 格回到原地。
  \end{center}
\end{frame}

\begin{frame}{算一算：\(8+7=\) ?}
  \imgph[4.2cm]{12 格圆盘：高亮 \(8\)，从 \(8\) 顺时针画一道弧，走过 7 格到 \(3\)}
  \pause
  \begin{center}
    {\large \(8+7=15\)，但走满一圈就回到原点，所以 \textbf{\(8+7\equiv 3 \pmod{12}\)}}
  \end{center}
\end{frame}

\begin{frame}{再算一个：\(11+1=\) ?}
  \imgph[4.2cm]{12 格圆盘：高亮 \(11\)，再走一格到 \(0\)（经过 12 点的位置）}
  \pause
  \begin{center}
    {\large \(11+1\equiv 0 \pmod{12}\)}\quad —— 从 11 再走一格，正好回到 0。
  \end{center}
\end{frame}

\begin{frame}{用手指回答}
  报出下面每个式子的结果（模 12）：
  \begin{itemize}
    \item \(10+5\)
    \item \(7+7\)
    \item \(9+4\)
    \item \(6+6\)
    \item \(3+3+3+3+3\)
  \end{itemize}
  \pause
  \begin{block}{答案}
    \(3\)；\quad \(2\)；\quad \(1\)；\quad \(0\)；\quad \(3\)
  \end{block}
\end{frame}

% ============================================================================
\section{单位元}

\begin{frame}{哪个数加进去，等于没加？}
  在模 12 的加法里，哪一个数 \(e\) 满足：
  \[
    a + e = a \quad\text{对每个 } a \text{ 都成立}
  \]
  \pause
  \begin{center}
    {\large 是 \textbf{0}。}\quad 因为 \(a+0 = a\)，走 0 格就是不动。
  \end{center}
\end{frame}

\begin{frame}{给它一个正式的名字}
  \begin{definition}[单位元]
    设 \(*\) 是集合 \(S\) 上的一个运算。如果存在 \(e \in S\)，使得
    \[
      a * e = e * a = a
    \]
    对集合里\textbf{每一个}元素 \(a\) 都成立，就称 \(e\) 是这种运算的\textbf{单位元}。
  \end{definition}
  \pause
  \begin{att}
    「单位」是「一个」的意思：它不改变任何东西。
  \end{att}
\end{frame}

% ============================================================================
\section{逆元}

\begin{frame}{每个数要走几格，才能回到 0？}
  把 12 个数两两配起来：
  \pause
  \begin{center}
  \begin{tabular}{@{}l*{12}{c}@{}}
    \toprule
    \(a\) & 0 & 1 & 2 & 3 & 4 & 5 & 6 & 7 & 8 & 9 & 10 & 11 \\
    \midrule
    \(\text{搭档}\) & 0 & 11 & 10 & 9 & 8 & 7 & 6 & 5 & 4 & 3 & 2 & 1 \\
    \bottomrule
  \end{tabular}
  \end{center}
  \pause
  \(1\) 配 \(11\)（\(1+11=12\equiv 0\)），\(2\) 配 \(10\)，\(3\) 配 \(9\)，\(\dots\)
  \pause
  \begin{att}
    有两个人只配自己：\textbf{\(0\)}（\(0+0=0\)）和 \textbf{\(6\)}（\(6+6=12\equiv 0\)）。
  \end{att}
\end{frame}

\begin{frame}{正式的名字：逆元}
  \begin{definition}[逆元]
    设 \(e\) 是运算 \(*\) 的单位元。如果对元素 \(a\) 能找到一个元素 \(b\)，使得
    \[
      a * b = b * a = e
    \]
    就称 \(b\) 是 \(a\) 的\textbf{逆元}，也说 \(a\) 与 \(b\) \textbf{互逆}。
  \end{definition}
  \pause
  \begin{block}{模 12 加法里}
    每个数的逆元都存在，就是它的「搭档」。
  \end{block}
\end{frame}

% ============================================================================
\section{和有理数对上}

\begin{frame}{这两样东西，你们早就在用了}
  \begin{center}
  \begin{tabular}{@{}lcc@{}}
    \toprule
     & \textbf{单位元} & \textbf{逆元} \\
    \midrule
    有理数加法 & \(0\) & 相反数 \(-a\) \\
    有理数乘法 & \(1\) & 倒数 \(\dfrac{1}{a}\ (a\neq 0)\) \\
    模 12 加法 & \(0\) & 搭档（\(a\) 配 \(12-a\)） \\
    \bottomrule
  \end{tabular}
  \end{center}
  \pause
  \begin{att}
    只是以前没这么叫过。\pause 负数是为了「加法能走回头路」才被造出来的：
    \(5\) 加上 \(-5\) 得到单位元 \(0\)。
  \end{att}
\end{frame}

\begin{frame}{这节课要带走的两句话}
  \begin{block}{一}
    取余数是有名字的算术：模 \(n\) 的加法。
  \end{block}
  \begin{block}{二}
    每一套这样的算术里，都有\textbf{单位元}和\textbf{逆元}——
    有理数加法、有理数乘法、模 12 加法，都有。
  \end{block}
  \pause
  \begin{center}
    {\large\textcolor{textgreen}{\bfseries 下节课换个运算试试：乘法还能不能找到逆元？}}
  \end{center}
\end{frame}

% ---------------- 结尾页 ----------------
\begin{frame}
  \begin{center}
    \vspace{1.6cm}
    {\Large\textcolor{textgreen}{\bfseries 下课}}

    \vspace{1.0cm}
    {\large\textcolor{accentgreen}{下节课：撤回一步——逆元}}
  \end{center}
\end{frame}

\end{document}